\documentclass[a4paper,11pt]{article}
\usepackage[T2A]{fontenc}
\usepackage[utf8]{inputenc}
\usepackage[russian]{babel}
\usepackage{amsmath,amssymb,mathtools}
\usepackage{geometry}
\geometry{left=22mm,right=22mm,top=21mm,bottom=21mm,headheight=22pt}
\usepackage{microtype}
\usepackage{array,booktabs,tabularx}
\usepackage{enumitem}
\usepackage{tikz}
\usetikzlibrary{arrows.meta,positioning,shapes.geometric}
\usepackage{fancyhdr}
\usepackage{setspace}
\usepackage{xcolor}
\usepackage{needspace}
\usepackage{hyperref}
\definecolor{accent}{HTML}{087D83}
\definecolor{darkaccent}{HTML}{075A61}
\definecolor{ink}{HTML}{30383A}
\definecolor{soft}{HTML}{667376}
\hypersetup{colorlinks=true,linkcolor=darkaccent,urlcolor=accent,pdfauthor={Лёвин Артём Александрович},pdftitle={Чек-лист — Схемы погашения кредита — Анастасия Павлова}}
\setstretch{1.14}
\setlength{\parindent}{0pt}
\setlength{\parskip}{5pt}
\setlist[itemize]{leftmargin=1.5em,topsep=3pt,itemsep=3pt,parsep=0pt}
\pagestyle{fancy}
\fancyhf{}
\fancyhead[L]{\small\color{soft} Схемы погашения кредита}
\fancyhead[R]{\small\color{soft} Анастасия Павлова}
\fancyfoot[C]{\small\color{soft}\thepage}
\renewcommand{\headrulewidth}{0.35pt}
\renewcommand{\headrule}{\hbox to\headwidth{\color{accent}\leaders\hrule height \headrulewidth\hfill}}
\newcommand{\topic}[1]{\needspace{8\baselineskip}\vspace{6pt}{\Large\color{darkaccent}\bfseries #1}\par\vspace{3pt}}
\newcommand{\subtopic}[1]{\needspace{5\baselineskip}\vspace{5pt}{\large\color{darkaccent}\bfseries #1}\par\vspace{1pt}}
\newcommand{\key}[1]{\textcolor{darkaccent}{\textbf{#1}}}
\newcommand{\ru}[1]{\text{#1}}
\begin{document}
\thispagestyle{empty}
\begin{center}
\vspace*{19mm}
{\large\color{accent} МАТЕМАТИКА}\par
\vspace{25mm}
{\Huge\color{ink} Чек-лист}\par
\vspace{13mm}
{\LARGE\color{accent} Схемы погашения кредита\\[3mm] и расчёт переплаты}\par
\vspace{13mm}
{\normalsize Произвольные, дифференцированные и аннуитетные платежи\\[2mm]
Разложение чисел на простые множители}\par
\vspace{23mm}
{\normalsize Подготовка к ЕГЭ по профильной математике}\par
\vfill
{\large Анастасия Павлова}\par
\vspace{11mm}
{9 ноября 2026 г.}\par
\vspace{25mm}
{\small Лёвин Артём Александрович\\[1mm]
эксклюзивно для Анастасии Павловой}\par
\vspace{17mm}
\begin{tikzpicture}
\draw[accent,very thin] (-3.1,0) .. controls (-2.1,0.63) and (-1.25,-0.42) .. (0,0)
 .. controls (1.25,0.47) and (2.05,-0.58) .. (3.1,0);
\end{tikzpicture}
\end{center}
\clearpage

\topic{1. Общая схема кредита}

Обозначим через $S$ первоначальную сумму кредита, через $n$ --- число лет, через
$p$ --- годовую процентную ставку \emph{в долях единицы}. Коэффициент увеличения долга
за один год равен
\[
 p=\frac{r}{100}\quad\text{при ставке }r\%,\qquad k=1+p.
\]
Например, $12\%=0{,}12$, поэтому $k=1{,}12$.

Для каждого года заполняем пять граф. Именно эта таблица работает при любой
схеме погашения.

\begin{center}
\small
\renewcommand{\arraystretch}{1.4}
\begin{tabularx}{\linewidth}{>{\raggedright\arraybackslash}X >{\centering\arraybackslash}X >{\centering\arraybackslash}X >{\centering\arraybackslash}X >{\centering\arraybackslash}X}
\toprule
Начальная сумма & Начисленные проценты & Общая задолженность & Платёж & Остаток долга \\
\midrule
$B_{i-1}$ & $pB_{i-1}$ & $kB_{i-1}$ & $P_i$ & $B_i$ \\
\bottomrule
\end{tabularx}
\end{center}

Здесь $B_{i-1}$ --- долг на начало года, $P_i$ --- платёж за год, $B_i$ --- долг
после внесения платежа. Главные равенства:
\[
 \begin{aligned}
 \ru{Общая задолженность}&=\ru{начальный долг}+\ru{проценты},\\
 B_i&=kB_{i-1}-P_i,\\
 P_i&=kB_{i-1}-B_i.
 \end{aligned}
\]

\key{Проверка полного погашения:} после последнего платежа $B_n=0$.

\subtopic{Общая сумма платежей и переплата}
\[
 T=P_1+P_2+\cdots+P_n,\qquad \Delta=T-S.
\]
Если дополнительных комиссий нет и кредит полностью погашен, переплата
равна сумме начисленных процентов:
\[
 \Delta=I_1+I_2+\cdots+I_n.
\]

\topic{2. Как различать условия}

\begin{itemize}
\item \key{Равные платежи каждый год:} аннуитетная схема; величина $P_i=x$ постоянна.
\item \key{Основной долг погашается равными частями:} дифференцированная схема; ежегодно возвращается $S/n$ основного долга, проценты постепенно уменьшаются.
\item \key{Задан иной закон изменения остатка или платежей:} составляем общую таблицу, следуя непосредственно условию.
\end{itemize}

Слова \emph{«равные платежи»} и \emph{«равные части основного долга»}
задают разные правила. Если способ погашения не указан, его нельзя
угадывать: используйте только данные условия.
\clearpage

\begin{center}
{\Large\color{darkaccent}\bfseries Алгоритм выбора схемы погашения}\par
\vspace{17mm}
\begin{tikzpicture}[>=Latex, node distance=13mm and 17mm,
 terminator/.style={ellipse,draw=darkaccent,thin,align=center,text width=51mm,minimum height=14mm,inner sep=4pt,font=\small},
 decision/.style={diamond,draw=accent,thin,aspect=2.4,align=center,text width=39mm,inner sep=2pt,font=\small},
 outcome/.style={rounded corners=2pt,draw=accent,thin,align=left,text width=57mm,minimum height=18mm,inner sep=6pt,font=\small},
 line/.style={->,draw=darkaccent,thin}]
\node[terminator] (start) {Прочитайте условие кредита};
\node[decision,below=17mm of start] (eq) {Все годовые платежи равны?};
\node[outcome,right=19mm of eq] (ann) {\textbf{Аннуитетная}\par Платёж $x$ постоянен,\par $B_i=kB_{i-1}-x$.};
\node[decision,below=19mm of eq] (part) {Основной долг гасят равными частями?};
\node[outcome,right=19mm of part] (diff) {\textbf{Дифференцированная}\par Погашение тела: $S/n$.\par Платёж уменьшается.};
\node[decision,below=19mm of part] (law) {Есть особый закон изменения остатков или платежей?};
\node[outcome,right=19mm of law] (custom) {\textbf{Произвольная}\par Сначала находим остатки,\par затем каждый платёж.};
\node[outcome,below=20mm of law] (unknown) {\textbf{Проверьте условие}\par Не делайте вывод о схеме без достаточных данных.};
\draw[line] (start) -- (eq);
\draw[line] (eq.east) -- node[above,font=\scriptsize] {да} (ann.west);
\draw[line] (eq.south) -- node[right,font=\scriptsize] {нет} (part.north);
\draw[line] (part.east) -- node[above,font=\scriptsize] {да} (diff.west);
\draw[line] (part.south) -- node[right,font=\scriptsize] {нет} (law.north);
\draw[line] (law.east) -- node[above,font=\scriptsize] {да} (custom.west);
\draw[line] (law.south) -- node[right,font=\scriptsize] {нет} (unknown.north);
\end{tikzpicture}
\end{center}
\clearpage

\topic{3. Произвольная схема: остаток уменьшается на 20\%}

Если по условию \emph{в конце каждого из первых двух лет} долг составляет
$80\%$ долга на начало соответствующего года, то
\[
 B_1=0{,}8S,\qquad B_2=0{,}8B_1=0{,}64S.
\]
На каждом шаге платёж равен общей задолженности за вычетом заданного остатка:
\[
 \begin{aligned}
 P_1&=kS-0{,}8S,\\
 P_2&=k\cdot0{,}8S-0{,}64S.
 \end{aligned}
\]
\key{Критически важно:} $80\%$ берём от долга \emph{на начало года}, а не от суммы
после начисления процентов. В последний год, если требуется полностью закрыть
кредит, остаток равен нулю.

\subtopic{Разобранный случай: 750 тыс. рублей, 12\% в год}
Кредит взят на три года. После первого и второго годов основной долг
сокращается на $20\%$; по окончании третьего года долг погашается полностью.
Все денежные суммы в таблице --- \textbf{в тыс. руб.}

\begin{center}
\small\renewcommand{\arraystretch}{1.42}
\begin{tabularx}{\linewidth}{c *{5}{>{\centering\arraybackslash}X}}
\toprule
Год & В начале & Проценты & Задолженность & Платёж & Остаток \\
\midrule
1 & 750 & 90 & 840 & 240 & 600 \\
2 & 600 & 72 & 672 & 192 & 480 \\
3 & 480 & $57{,}6$ & $537{,}6$ & $537{,}6$ & 0 \\
\bottomrule
\end{tabularx}
\end{center}

Первый платёж: $840-600=240$. Второй: $672-480=192$.
В третьем году выплачивается вся задолженность: $480\cdot1{,}12=537{,}6$.
\[
 T=240+192+537{,}6=969{,}6\ \text{тыс. руб.}
\]
Переплата составит $969{,}6-750=219{,}6$ тыс. руб.

\topic{4. Дифференцированная схема}

Основной долг делится на $n$ одинаковых частей. В каждом году
погашается $S/n$ основного долга и оплачиваются проценты на непогашенную часть.
Платёж зависит от номера года $i$:
\[
 \boxed{P_i=\frac{S}{n}+pS\,\frac{n-i+1}{n}},\qquad 1\le i\le n.
\]
Первый платёж самый большой; последний --- самый маленький.
\clearpage

\topic{4. Дифференцированная схема: пять формул}

Все пять выражений относятся к кредиту на $n$ лет с равными годовыми долями
погашения основного долга. Ставка $p$ записана в долях единицы.

\begin{align*}
\textbf{1. }\quad T&=S+\frac{Sp(n+1)}{2}
 &&\text{общая сумма платежей},\\[4pt]
\textbf{2. }\quad \Delta&=\frac{Sp(n+1)}{2}
 &&\text{переплата в рублях},\\[4pt]
\textbf{3. }\quad q&=\frac{p(n+1)}{2}\cdot100\%
 &&\text{процент переплаты},\\[4pt]
\textbf{4. }\quad P_1&=Sp+\frac{S}{n}
 &&\text{первый (наибольший) платёж},\\[4pt]
\textbf{5. }\quad P_n&=\frac{Sp}{n}+\frac{S}{n}
 &&\text{последний (наименьший) платёж}.
\end{align*}

Формула переплаты получается из первой вычитанием $S$; процент переплаты
получается делением $\Delta$ на $S$ и умножением на $100\%$.

\subtopic{Разобранный случай: 600 тыс. рублей, 8\%, три года}
Каждый год возвращается $600/3=200$ тыс. руб. основного долга.

\begin{center}
\small\renewcommand{\arraystretch}{1.42}
\begin{tabularx}{\linewidth}{c *{5}{>{\centering\arraybackslash}X}}
\toprule
Год & В начале & Проценты & Задолженность & Платёж & Остаток \\
\midrule
1 & 600 & 48 & 648 & 248 & 400 \\
2 & 400 & 32 & 432 & 232 & 200 \\
3 & 200 & 16 & 216 & 216 & 0 \\
\bottomrule
\end{tabularx}
\end{center}

\[
 \begin{aligned}
 T&=248+232+216=696\ \text{тыс. руб.},\\
 \Delta&=48+32+16=96\ \text{тыс. руб.},\\
 q&=\frac{96}{600}\cdot100\%=16\%.
 \end{aligned}
\]
Полученные суммы совпадают с формулами: $600+600\cdot0{,}08\cdot(3+1)/2=696$.

\key{Проверка:} проценты снижаются вместе с остатком долга, поэтому
$248>232>216$.
\clearpage

\topic{5. Аннуитетная схема: равные платежи}

При аннуитете все годовые платежи равны одной величине $x$. Проценты
начисляются на остаток долга, поэтому таблица строится по рекуррентному правилу
\[
 B_i=kB_{i-1}-x,\qquad B_0=S,\quad B_n=0.
\]
Для трёх лет с годовой ставкой $20\%$ имеем $k=1{,}2$:
\[
 \begin{aligned}
 B_1&=1{,}2S-x,\\
 B_2&=1{,}2(1{,}2S-x)-x
      =1{,}44S-2{,}2x,\\
 B_3&=1{,}2(1{,}44S-2{,}2x)-x
      =1{,}728S-3{,}64x.
 \end{aligned}
\]
Поскольку после третьего платежа долг погашен,
\[
 1{,}728S-3{,}64x=0,\qquad
 x=\frac{1{,}728}{3{,}64}S=\frac{216}{455}S.
\]

\subtopic{Как правильно прочитать условие о переплате}
Если \emph{общая сумма выплат на $96\,500$ руб. больше суммы, взятой в кредит},
то разность равна $96\,500$:
\[
 \underbrace{3x}_{\text{все платежи}}-\underbrace{S}_{\text{взяли в банке}}=96\,500.
\]
Подставим найденный платёж:
\[
 3\cdot\frac{216}{455}S-S
 =\frac{648-455}{455}S
 =\frac{193}{455}S=96\,500.
\]
Отсюда
\[
 S=96\,500\cdot\frac{455}{193}=227\,500\ \text{руб.}
\]
Величина ежегодного платежа и итоговая сумма:
\[
 x=\frac{216}{455}\cdot227\,500=108\,000\ \text{руб.},
 \qquad T=3x=324\,000\ \text{руб.}
\]

\key{Контроль вычислений:} $324\,000-227\,500=96\,500$ руб.
\clearpage

\topic{5. Аннуитет: проверка по таблице}

Подставим $S=227\,500$ руб. и $x=108\,000$ руб. в таблицу. Результат
должен дать нулевой остаток после третьего года.

\begin{center}
\small\renewcommand{\arraystretch}{1.42}
\begin{tabularx}{\linewidth}{c *{5}{>{\centering\arraybackslash}X}}
\toprule
Год & В начале & Проценты & Задолженность & Платёж & Остаток \\
\midrule
1 & 227\,500 & 45\,500 & 273\,000 & 108\,000 & 165\,000 \\
2 & 165\,000 & 33\,000 & 198\,000 & 108\,000 & 90\,000 \\
3 & 90\,000 & 18\,000 & 108\,000 & 108\,000 & 0 \\
\bottomrule
\end{tabularx}
\end{center}

Сумма начисленных процентов составляет
$45\,500+33\,000+18\,000=96\,500$ руб. Это независимая проверка переплаты.

\topic{6. Разложение на простые множители}

\key{Простое число} имеет ровно два положительных делителя: единицу и само число.
Любое целое число, большее единицы, можно единственным образом разложить
на простые множители (если не учитывать порядок множителей).

При вычислениях с аннуитетом возникают дроби с большими числами.
Удобнее сперва разложить их:
\[
 \begin{aligned}
 1728&=2^6\cdot3^3,\\
 3640&=2^3\cdot5\cdot7\cdot13,\\
 96\,500&=193\cdot2^2\cdot5^3.
 \end{aligned}
\]
Число $193$ простое: оно не делится на $2,3,5,7,11,13$; проверять
большие простые делители не требуется, так как $\sqrt{193}<14$.

Поэтому
\[
 \frac{1728}{3640}
 =\frac{2^6\cdot3^3}{2^3\cdot5\cdot7\cdot13}
 =\frac{216}{455},\qquad
 \frac{96\,500}{193}=500.
\]

\key{Практическое правило:} в больших дробях сначала сокращайте общие
простые множители, затем перемножайте оставшиеся числа.

\topic{7. Ошибки, которых важно избегать}
\begin{itemize}
\item При вычислении $20\%$ от остатка берите правильную базу: долг
\emph{на начало года}, если именно так сказано в условии.
\item В дифференцированной схеме фиксирована доля \emph{основного долга},
поэтому вся сумма платежа каждый год меняется.
\item В аннуитетной схеме фиксирован \emph{полный платёж}, а сумма процентов
и погашения основного долга внутри него изменяется.
\item После раскрытия скобок складывайте коэффициенты внимательно:
$2{,}64x+x=3{,}64x$.
\item Фраза «выплачено на $96\,500$ руб. больше» означает
$T-S=96\,500$, тогда как «выплачено всего $96\,500$ руб.» означала бы
$T=96\,500$.
\item Сумма платежей и переплата имеют размерность \emph{рубли}; процент
переплаты имеет размерность \emph{проценты}.
\end{itemize}
\clearpage

\topic{Итоговая памятка}
\small
\renewcommand{\arraystretch}{1.53}
\begin{tabularx}{\linewidth}{>{\raggedright\arraybackslash}p{0.35\linewidth} >{\raggedright\arraybackslash}X}
\toprule
Сигнал в условии & Действие \\
\midrule
Любая кредитная задача & Начальная сумма $\to$ проценты $\to$ общая задолженность $\to$ платёж $\to$ остаток.\\
Равные полные платежи & Обозначить каждый платёж $x$; использовать $B_i=kB_{i-1}-x$.\\
Равные доли основного долга & Выделить $S/n$ ежегодно; прибавлять проценты на текущий остаток.\\
Особый закон изменения долга & Сначала определить остаток, затем платёж как разность задолженности и остатка.\\
Просят сумму выплат & Сложить все платежи: $T=\sum P_i$.\\
Просят переплату & Вычесть исходную сумму: $\Delta=T-S$; проверить сумму процентов.\\
Большая дробь & Разложить числитель и знаменатель на простые множители и сократить.\\
\bottomrule
\end{tabularx}
\normalsize

\vspace{9mm}
\begin{center}
\color{darkaccent}\large
\textbf{Один контрольный принцип для всех схем:}
\end{center}
\[
 \text{остаток после платежа}
 =\text{долг с процентами}-\text{платёж}.
\]

Для полного погашения последняя строка таблицы заканчивается нулём.
Эта проверка помогает обнаружить ошибки независимо от выбранной схемы.

\vfill
\begin{center}
{\small\color{soft} Лёвин Артём Александрович --- эксклюзивно для Анастасии Павловой}
\end{center}
\end{document}
